\documentclass[main.tex]{subfiles}

% Individual lecture title
\ifSubfilesClassLoaded{%
    \title{Math 113: Irreducible polynomial for alpha over F}
    \author{\defaultauthor}
    \date{November 19, 2026}
}{}

\begin{document}

\ifSubfilesClassLoaded{\maketitle}{}

\section{Lecture 24: Irreducible polynomial for \texorpdfstring{$\alpha$}{α} over F}

An element $\alpha \in E  \supseteq F$ can be a zero for many polynomials $f(x) \in F[x]$, but: 

\begin{theorem}
    Let $\alpha \in E \supseteq F$ be algebraic over F. Then there is an irreducibel $p(x) \in F[x]$ with $p(\alpha) = 0$

    This p(x) is unique up to a constant facotr in F and is a polynomial of minimal degree in F[x] having $\alpha$ as a zero. 
\end{theorem}
\begin{definition}
    This polynomial is the \textbf{irreducible polynomial for } $\alpha$ over F. We denote it by $irr(\alpha,F)$

    The degree of $irr(\alpha,F)$ is the degree of $\alpha$ over F, denoted $deg(\alpha,F)$
\end{definition}
\begin{example}
    $irr(\sqrt{2},\mathbb{Q})=x^2 - 2$, so $deg(\sqrt{2}, \mathbb{Q}) = 2 $
\end{example}
\begin{example}
    $irr(\sqrt{1+\sqrt{3}}, \mathbb{Q} = x^4 - 2x^2-2)$, so $deg(\sqrt{1+\sqrt{3}},\mathbb{Q})=4$
\end{example}
\subsection{Adjoining elements to Fields}
Let E be an extension field of F, $\alpha \in E$. So 
$$\varphi_{\alpha}: F[x] \to E$$ 

We want to define 
\end{document}